% © 2026 Ing. David Jaroš — CC BY-NC-ND 4.0
%
% This work is licensed under a Creative Commons Attribution-NonCommercial-NoDerivatives
% 4.0 International License (CC BY-NC-ND 4.0).
%
% License History: Earlier drafts (up to v0.3) were released under CC BY 4.0.
% From v0.4 onward, all material is released under CC BY-NC-ND 4.0 to protect
% the integrity of the theoretical work during ongoing academic development.
%
% See LICENSE.md for full license text.

\ifdefined\INCLUDEMODE
\else
  \documentclass[12pt,a4paper]{article}
  \usepackage[a4paper, margin=2.5cm]{geometry}
  \usepackage{amsmath, amssymb, amsthm}
  \usepackage{mathtools}
  \usepackage{hyperref}
  \usepackage{xcolor}
  \usepackage{mdframed}
  \usepackage{booktabs}

  \newtheorem{theorem}{Theorem}[section]
  \newtheorem{lemma}[theorem]{Lemma}
  \newtheorem{proposition}[theorem]{Proposition}
  \theoremstyle{definition}
  \newtheorem{definition}[theorem]{Definition}
  \theoremstyle{remark}
  \newtheorem{remark}[theorem]{Remark}

  \newmdenv[backgroundcolor=yellow!20, linecolor=orange, linewidth=1.5pt]{gapbox}
  \newmdenv[backgroundcolor=green!15, linecolor=green!60!black, linewidth=1.5pt]{resultbox}
  \newmdenv[backgroundcolor=red!10, linecolor=red!60, linewidth=1.5pt]{deadendbox}
  \newmdenv[backgroundcolor=blue!8, linecolor=blue!50, linewidth=1.5pt]{insightbox}

  \title{\textbf{Electron Mass --- Step 2: Geometric $R_\psi$ Approach}\\
         \large Can the Zero-Mode Energy of $S[\Theta]$ Fix $m_e$?}
  \author{David Jaro\v{s}}
  \date{2026-03-06}

  \begin{document}
  \maketitle

  \begin{abstract}
  We investigate whether the $\psi$-circle radius $R_\psi$ can be fixed
  non-circularly from the kinetic action $S[\Theta]$.  Three candidates
  for fixing $R_\psi$ (self-duality, winding consistency, modularity) have
  previously been shown to fail (\texttt{docs/ALPHA\_FROM\_ME\_ANALYSIS.md}).
  We revisit the zero-mode energy approach:
  $E_0 = \hbar/(2R_\psi) = m_e c^2$ gives $R_\psi = \hbar/(2m_e c)$,
  differing from the standard calibration $R_\psi = \hbar/(m_e c)$
  by a factor of~2.  We test whether this factor-of-2 variant is
  consistent with other UBT constraints.  The approach is a dead end:
  $R_\psi$ cannot be fixed independently of $m_e$ within current UBT.
  \end{abstract}
\fi

%──────────────────────────────────────────────────────────────────────────────
\section{The Zero-Mode Energy Approach}
\label{sec:zero_mode}
%──────────────────────────────────────────────────────────────────────────────

On a circle of radius $R$, the quantum zero-point energy of a scalar
field is (standard quantum mechanics):
\begin{equation}
  E_0 \;=\; \frac{\hbar}{2R_\psi}
  \label{eq:zero_mode}
\end{equation}
in natural units.  If we identify this with the electron rest-mass energy:
\begin{equation}
  \frac{\hbar}{2R_\psi} \;=\; m_e c^2
  \quad\Longrightarrow\quad
  R_\psi \;=\; \frac{\hbar}{2 m_e c} \;=\; \frac{\bar\lambda_e}{2},
\end{equation}
where $\bar\lambda_e = \hbar/(m_e c)$ is the reduced Compton wavelength.
The current UBT calibration uses $R_\psi = \bar\lambda_e$ (without the
factor of~2).

\subsection{Comparison with Current Calibration}

\begin{center}
\begin{tabular}{ll}
  \toprule
  Formula & $R_\psi$ \\
  \midrule
  Zero-mode energy $E_0 = m_e c^2$ & $\hbar/(2m_e c) = \bar\lambda_e/2$ \\
  Current UBT calibration & $\hbar/(m_e c) = \bar\lambda_e$ \\
  \bottomrule
\end{tabular}
\end{center}

The factor-of-2 difference changes the KK mode spectrum:
$M_n = n/(2R_\psi) = n m_e c/\hbar$ vs $M_n = n m_e c/\hbar$ (no factor).
This affects the KK tower in the vacuum polarisation but \emph{not}
the one-loop $\beta$-function coefficient $B_0 = 8\pi$ (which is
independent of $R_\psi$ at one loop).

%──────────────────────────────────────────────────────────────────────────────
\section{Consistency Test}
\label{sec:consistency}
%──────────────────────────────────────────────────────────────────────────────

If $R_\psi = \hbar/(2m_e c)$, the winding-number excitation energies are:
\begin{equation}
  E_n \;=\; \frac{n^2\hbar^2}{2M_\Theta R_\psi^2}
       \;=\; \frac{n^2\hbar^2}{2M_\Theta} \times \frac{4m_e^2 c^2}{\hbar^2}
       \;=\; 2n^2 M_\Theta m_e^2 c^4 / \hbar^2.
\end{equation}
For $M_\Theta = m_e$: $E_n = 2n^2 m_e c^2$.  For $n=1$: $E_1 = 2m_e c^2$,
twice the electron rest mass.  This would identify the $n=1$ mode with a
particle of mass $m_{\Theta,1} = \sqrt{2}\,m_e$, not $m_e$.

This is inconsistent with the identification $\Theta_1 \leftrightarrow$
electron (generation 1) at the electron mass scale.

\subsection{Previous Dead Ends}
From \texttt{canonical/geometry/biquaternionic\_vacuum\_solutions.tex~§2}
and \texttt{docs/ALPHA\_FROM\_ME\_ANALYSIS.md~§6.2}:
\begin{itemize}
  \item \textbf{Self-duality condition}: $R_\psi = \ell_P$ (Planck length)
        --- gives $m_e \sim m_P$, wrong by 22 orders of magnitude.
  \item \textbf{Winding consistency}: $R_\psi = n^*/(\Lambda_{\mathrm{UV}})$
        --- circular (depends on $n^* = 137$).
  \item \textbf{Modularity condition}: no modular invariant condition
        fixes $R_\psi$ at $\bar\lambda_e$ from first principles.
\end{itemize}

%──────────────────────────────────────────────────────────────────────────────
\section{Conclusion}
\label{sec:conclusion}
%──────────────────────────────────────────────────────────────────────────────

\begin{deadendbox}
\textbf{Dead end (geometric $R_\psi$ approach):}

None of the geometric conditions tested fixes $R_\psi = \hbar/(m_e c)$
or $R_\psi = \hbar/(2m_e c)$ from first principles:
\begin{itemize}
  \item The zero-mode energy gives $R_\psi = \hbar/(2m_e c)$, which is
        \emph{inconsistent} with the electron-mass identification.
  \item The standard calibration $R_\psi = \hbar/(m_e c)$ is circular:
        it uses $m_e$ as input to define $R_\psi$.
  \item Three prior attempts (self-duality, winding, modularity) failed
        (documented in \texttt{docs/ALPHA\_FROM\_ME\_ANALYSIS.md}).
\end{itemize}

\textbf{Honest statement:}
Gap~Y2 (derive $v$ independently of $m_e$) is not closeable by the
geometric $R_\psi$ approach within current UBT.  $R_\psi$ remains a
\emph{calibrated} parameter.

\textbf{Next step:} Approach the problem from modular forms
(Step~3 of this series).
\end{deadendbox}

\ifdefined\INCLUDEMODE
\else
  \end{document}
\fi
